\chapter{Conclusion and Future Work \label{ch:conclusion}}
\chaptertoc

This chapter concludes the thesis by synthesizing findings on graph-based representation learning for architecture recovery and C2A mapping. It summarizes the conclusions across the three research objectives and outlines directions for further research arising from the findings and limitations of the work.

\section{Conclusion \label{sec:conclusion}}

This thesis examined how graph-based representation learning can support software architecture recovery and C2A mapping, and what different implementation signals contribute to these tasks. Recovery was studied as an unsupervised task and mapping as a supervised task with limited labeled data. Across five studies, the work progressed from learning joint representations of static dependencies, folder organization, and source-code content to examining their individual contributions and their interactions with graph structure.

The first objective examined whether dependency, lexical, and organizational signals could be combined in learned representations for recovery and mapping, and how performance depended on the learning objective and model configuration. The joint representations supported both tasks, although the evidence differed between them. In recovery, the approaches of Papers~I and~II improved on or matched the evaluated baselines across most comparisons. In mapping, however, improvement over the Naive Bayes baseline was not specific to the GNN models. Both the GNN models of Paper~III and an MLP using the same features outperformed that baseline. This improvement therefore does not by itself establish a benefit from message passing.

In the recovery comparison using a common GAT encoder and input features, the reconstruction-based approach generally produced better agreement with the reference decompositions than the contrastive approach. The approaches retained different training settings, so the comparison does not isolate the learning objective alone. Encoder choice also mattered: GAT generally outperformed GCN in both tasks. Modeling dependency relations separately improved mapping with lexical and folder-token features (Paper~III), although the encoder also had more parameters. With file vectors without folder-path features, the benefit was smaller and varied across representations and systems, as shown in the cross-study mapping evaluation. Separate relations gave no benefit in preliminary tests in dependency-only recovery (Paper~IV). Relation-specific modeling was therefore useful under some input representations and tasks, rather than consistently beneficial.

The joint-signal comparisons did not separate the contributions of the individual signals. The second objective therefore examined what architectural information static dependencies and source-code content carry, and how representation choices make that information available. In recovery from static dependencies alone, the results suggest that the extent of dependency context encoded in a file's representation (structural reach) matters. Node2Vec extends this reach through random walks and GAT-DC through PPR diffusion. Both outperformed the classical baselines on most systems, whereas the two-layer GAT, limited to its two-hop neighborhood, was generally weaker. Agreement with the reference decomposition was lower where within-module and cross-module dependencies were similarly dense, whereas the dependency cohesion of the recovered partitions frequently matched or exceeded that of the reference. Dependency structure alone therefore supported recovery to a bounded extent: its usefulness depended on structural reach, on how well the reference modules were separated by dependencies, and on the evaluation criterion.

In mapping with an MLP, Word2Vec and zero-shot UniXcoder showed no consistent ordering, so the pretrained code transformer was not consistently better than a lexical representation. Dependency-guided fine-tuning improved UniXcoder on nearly all systems, supporting adaptation of the pretrained representation to the target system using dependency information. Node2Vec and fine-tuned UniXcoder generally outperformed the other two representations. Neither of the two was uniformly best: Node2Vec most often performed best, while fine-tuned UniXcoder was significantly stronger on some systems. Dependency information was therefore useful in mapping as in recovery, regardless of representation family, although differences in content and training prevent attributing the performance differences to dependency information alone.

The information carried by individual file representations did not establish what additional graph context could contribute. The third objective therefore examined how dependency-based message passing, folder organization, and member-level detail changed their usefulness. With undirected aggregation and fixed file representations, GAT and an MLP supplied with mean neighbor vectors improved on most systems for every representation, with no significant losses. Gains were largest for zero-shot UniXcoder and smallest for Node2Vec. Fine-tuned UniXcoder gained more than Node2Vec and had the highest median macro-F1 across systems for every predictor. The benefit of aggregation therefore depended on how dependency context was already encoded.

Aggregation direction also mattered. Source-to-target aggregation lowered median macro-F1 with Node2Vec, and the preferred direction depended on the system rather than the predictor or representation. Papers~III and~V agreed on some systems but differed on others, so neither the benefit nor the preferred direction can be assumed to transfer across setups. Undirected aggregation also added reverse edges, so direction and neighborhood size were not separated.

Folder organization supplied architectural evidence beyond static dependencies alone. Across the recovery studies, methods using organizational information agreed more closely with the reference decomposition, particularly in Paper~IV where dependencies gave limited evidence of module boundaries, but produced less dependency-cohesive partitions than file-level dependency-only methods. Reference agreement and dependency cohesion therefore favored different configurations. In mapping with HGT and source-to-target dependency aggregation, folder nodes improved every file representation, but the model with fine-tuned UniXcoder did not consistently outperform the directory-based baseline, limiting the evidence for a benefit beyond directory organization itself (Paper~V). Folder information and the reference decomposition may also overlap because directories are consulted during reference construction; the extent of this overlap was not quantified.

Member-level detail gave no consistent improvement in the evaluated configurations. With HGT and source-to-target dependency aggregation, member nodes lowered median mapping F1 for all three representations without folders and had representation-dependent effects with folders (Paper~V). Since member nodes also changed how features entered the model, their effects cannot be attributed to the added structure alone. In the cross-study recovery evaluation, member-level edge attributes gave no significant improvement over the same model without these attributes. These findings therefore concern the evaluated encoders and representations, not the usefulness of member-level information in general.

Taken together, dependency information remained useful, but encoding it in file representations did not make message passing redundant, and folder organization supplied further information. The dependence on encoder choice observed under Objective~1 recurred as a dependence on predictor and aggregation direction, and across the studies no representation or configuration was uniformly preferable: the contribution of each signal depended on how it was represented and combined with other evidence.

\section{Future Work~\label{sec:future_work}}
A first direction is to determine which aspects of graph context account for the observed performance differences. Undirected aggregation changed both propagation direction and neighborhood size, while member nodes changed graph structure and feature placement. Comparing incoming, outgoing, and combined neighborhoods at matched sizes could help distinguish direction from neighborhood size, while member-level comparisons should preserve feature content. Parameter-matched comparisons of relation-specific and collapsed-relation models could distinguish relational detail from model capacity.
 
Paper~I included interaction frequency and type together as member-level edge attributes. The cross-study recovery evaluation found no significant improvement from these attributes, but did not isolate the contribution of frequency. Whether frequency adds architectural information, and how it should enter representation learning through edge weights, attention terms, or relation encodings, remains open for both recovery and mapping.
 
A second direction concerns configuration selection and evaluation on unfamiliar and evolving systems. The benefit and preferred direction of neighbor aggregation varied across systems, but dependency homophily, used to interpret these differences, requires reference module assignments. Future work should examine whether properties observable in the dependency graph alone can guide configuration selection, with selection rules evaluated on systems not used to develop them. Folder-aware methods could also be assessed across systems that differ in how closely directories match the reference modules. Reporting performance separately for files whose module differs from that of their directory would help establish whether the methods provide useful predictions beyond directory organization.
 
The mapping evaluations made all system files available, leaving performance on subsequently introduced files untested. Future work could evaluate whether a model trained on mapped files from one version can assign new and changed files in later versions without retraining. The version history could also supply evolutionary signals, such as files that change together, which may carry architectural information that dependencies, folders, and code content do not. Practitioner studies could then assess whether the predictions support architecture maintenance.
